% RESEARCH_DRAFT_SHA256: 4db86759781a31627fb2e45def35521b693948abc380faaea2214c8e3eefdb90
% Layout and source-reference wrappers expanded for the web renderer.
\documentclass[11pt]{article}
\usepackage[margin=25mm]{geometry}
\usepackage{fontspec,unicode-math}
\setmainfont{Latin Modern Roman}
\setmathfont{Latin Modern Math}
\usepackage{amsmath,amssymb,mathtools,xurl,hyperref}
\hypersetup{colorlinks=true,urlcolor=blue,linkcolor=blue}
\setlength{\parindent}{0pt}
\setlength{\parskip}{5pt}
\setlength{\emergencystretch}{3em}
\newcommand{\sref}[2]{\hyperlink{source:#1:#2}{[S#2]}}
\newcommand{\cataloguescope}[1]{\paragraph{Source question.}#1}
\newcommand{\floor}[1]{\left\lfloor#1\right\rfloor}
\begin{document}

\section*{Pivot arborescences: local choices and a complete normalized-weight polytope construction}
\noindent catalogue entry; catalogue entry. Research attempt dated 11 October 2026.
\paragraph{Result in scope.}This manuscript contains a source-specific attempt and complete proofs of the stated auxiliary calculations. It makes no new solution or novelty claim for the entire target.
\section{Original problem and preserved definition}
The following definition and source list are reproduced from the repository. Any earlier inspection dates and reported literature claims in this preserved text describe that source record, not additional verification by this attempt.
% BEGIN PRESERVED DEFINITION
\subsection{Definitions and mathematical statement}

\paragraph{Shared notation.}Standard number systems are denoted by $\mathbb N,\mathbb Z,\mathbb Q,\mathbb R,\mathbb C$; index ranges, computational models and specialized notation follow the cited problem source and the definitions below.


For a nondegenerate linear program on a polytope and a generic linear objective, orient each edge toward decreasing objective value. A memoryless pivot rule chooses an improving outgoing edge using the current basis and fixed problem data, without a history of previous pivots. Choosing one such edge at every nonoptimal vertex yields a directed spanning tree toward the unique optimum, called an arborescence.
\paragraph{Statement recorded in the supplied source.}
For a fixed linear-programming objective, how do the induced arborescences vary with the memoryless pivot rule, and for a fixed memoryless pivot rule, how do they vary with the objective? \href{https://doi.org/10.4171/owr/2021/53}{[S2]}
\paragraph{Scope qualification.}
Degenerate bases, objective ties, and rules that permit cycling require conventions beyond this generic description. The supplied question does not identify those conventions, so this context explains the usual arborescence setting rather than extending the claim to every LP presentation.
\paragraph{Source status.}
The supplied catalogue labels this record partially solved. Its source-status metadata was checked against the supplied snapshot on 2026-10-11. This import does not independently establish current open status or mathematical correctness. The catalogue literature assessment is dated 2026-08-22.
\subsection{Short English statement}
Understand the directed trees produced by pivot rules as the rule or the objective changes.
\subsection{Sources}
\begin{itemize}
\item[S1] Ulam AI and the UnsolvedMath Contributors, UnsolvedMath: A Curated Collection of Open Mathematics Problems, catalogue entry. Catalogue curation and metadata: CC BY 4.0. The supplied source question is retained; definitions and the English summary are editorial additions.
\url{https://huggingface.co/datasets/ulamai/UnsolvedMath}.
\item[S2] Combinatorial Optimization (2022), Oberwolfach Reports, DOI 10.4171/owr/2021/53, pp. 2893--2954
\url{https://doi.org/10.4171/owr/2021/53}.
\item[S3] Combinatorial Optimization, Oberwolfach Reports 53/2021, primary report PDF
\url{https://ems.press/content/serial-article-files/46931}.
\textit{Underlying Oberwolfach Reports material retains its source-specific licence, including CC BY-SA 4.0 where applicable; see THIRD\_PARTY\_NOTICES.md.}
\end{itemize}
% END PRESERVED DEFINITION
\subsection{Verbatim supplied catalogue question}
For a fixed linear-programming objective, how do the induced arborescences vary with the memoryless pivot rule, and for a fixed memoryless pivot rule, how do they vary with the objective?

\section{A precise subclass of the broad question}
The source studies normalized-weight rules as a controlled subclass
of memoryless rules. We rederive its elementary polyhedral mechanism
for a fixed objective and state exactly what changes when the
objective varies. This is a proof within that subclass, not a
classification of every unrestricted memoryless rule.

\paragraph{Independent choices on an acyclic oriented graph.}
Let a finite directed graph be acyclic and have a unique sink $z$.
At every other vertex $v$ choose one outgoing arc. Following
chosen arcs cannot cycle, so it terminates at a sink, necessarily
$z$. The chosen graph is connected after forgetting directions,
and it has exactly one arc for each of the $N-1$ nonsink
vertices. A connected undirected graph on $N$ vertices with
$N-1$ edges is a tree: if it had a cycle, deleting one cycle
edge would preserve connectivity, while a connected graph has at
least $N-1$ edges by growing a spanning tree. Therefore each
choice gives a rooted arborescence. Conversely such an arborescence
chooses exactly one outgoing arc at each nonsink. If $d^+(v)$
denotes outdegree, their number is
$\prod_{v\ne z}d^+(v)$. This concerns local choices with no
shared rule-parameter restriction. It does not imply that one
normalized-weight vector realizes every such combination.

\paragraph{A complete fixed-objective Minkowski-sum construction.}
Fix an edge orientation on a polytope with strictly improving
directions and the unique-sink property just described. At each
nonsink $v$, let $D_v$ contain the vectors
$(u-v)/\eta(u-v)$ for improving neighbors $u$, where the
normalization values are fixed and positive. Distinct neighboring
edges have distinct positive rays, so these normalized vectors
are distinct: otherwise one neighbor would lie in the segment
from $v$ to another, contrary to being an edge endpoint/vertex.
Put $Q_v=\operatorname{conv}D_v$ and
\[
\Pi=\sum_{v\ne z}Q_v,
\]
where the sum consists of all sums of one point from each polytope.
For a weight vector $w$ giving a unique maximum in each finite
$D_v$, the rule selects its maximizing neighbor at each vertex.
The unique point maximizing $w\cdot x$ on $\Pi$ is the
sum of those selected directions: each summand is bounded by its
own maximum, and equality in the sum forces equality in every
summand. This proves both existence and uniqueness of that maximizer.

Moreover two such generic weights give the same maximizing point
of $\Pi$ if and only if they give the same selected neighbor
at every vertex. The reverse implication is immediate. For the
forward direction write the common point as the sum of the choices
for the second weight and evaluate the first weight on it. The
sum attains the first weight's sum of individual maxima. Each
individual summand is no larger than its maximum, so every summand
attains it; uniqueness forces every selected direction to agree.
Distinctness of directions forces the same neighbor. Thus the
maximizing-vertex regions of this explicit polytope encode exactly
the coherent arborescences in this fixed-normalization subclass.

\paragraph{Varying the objective and a finite comparison arrangement.}
For a rule with $w=c$ and normalization independent of $c$,
the orientation of each edge changes only when $c\cdot(u-v)=0$.
For fixed $v$, two candidate scores exchange order only when
\[
c\cdot\left(\frac{u-v}{\eta(u-v)}-
             \frac{u'-v}{\eta(u'-v)}\right)=0.
\]
On each connected open region avoiding all these finitely many
hyperplanes, every edge sign and every score comparison remains
constant. Therefore the whole chosen arborescence remains constant.
This proves an exact finite structural description for this rule
class, including the inequalities to test a proposed region. It
does not bound the number of regions polynomially in a general
input description, and it requires the vertices, adjacency data and
normalization values. Enumerating that data can itself be expensive.

\section{Audit and scope}
The source's Definitions 1--2 and normalized-weight construction on
printed pages 2936--2938 were inspected. These are known-style
rederivations of the displayed source mechanism; no novelty is claimed.
The audit kept generic tie-free parameters explicit and did not
equate arbitrary independent local choices with a common coherent
weight vector. Rules with history, objective-dependent normalization,
degeneracy or unrestricted tie conventions need separate treatment.
No enumeration or geometric experiment was run. The broad variation
question remains unresolved beyond the precisely proved subclass.

\section{Process and contributions}
Mathematics and model-assisted adversarial self-audit: OpenAI gpt-6-astra, ultra reasoning, as dispatched for this research batch; exact serving revision not exposed. Codex (GPT-6; exact serving revision not exposed) materially assisted with orchestration and format review. No human mathematical authorship, mathematical edits, or independent human proof checking is asserted.
Contribution prompt A was applied to the full definition and supplied source row, then prompt B was applied adversarially to the derived results. The item-specific audit above records the scope checks and rejected extrapolations. Local assembly and structural checks do not establish mathematical correctness. The source catalogue status is preserved in TOP\_PROBLEM and is separate from this attempt status.
\paragraph{Integration review.} Codex (GPT-6; exact serving revision not exposed) reviewed the written arguments and expanded layout and source-reference wrappers. Repository metadata and the production web renderer were checked using Python 3.13 and JavaScript with MathJax 3.2.2. This model-assisted review is not independent human proof verification; no PDF compilation or formal proof checking is asserted.

\section{Completion Estimate}
\noindent\textbf{COMPLETION: 0\%}
The proved material is a rederivation of known or auxiliary facts rather than a substantiated new advance on the original target. The normalized-weight subclass has an explicit proved polyhedral/comparison description; no unrestricted classification or efficient enumeration for all memoryless rules is supplied.
\section{References}
The complete original Sources subsection above is retained, including catalogue attribution and licence notices. The additional list records primary sources actually consulted for this attempt on 11 October 2026; uncited source links reproduced above were not all rechecked.
\begin{itemize}
\item Oberwolfach Report 53/2021, Combinatorial Optimization, original question and surrounding discussion read from the primary PDF text on 11 October 2026. Historical statements are not a current-status survey. \url{https://ems.press/content/serial-article-files/46931}.
\end{itemize}
\section{Final status}
\textbf{Unresolved in this attempt.} The normalized-weight subclass has an explicit proved polyhedral/comparison description; no unrestricted classification or efficient enumeration for all memoryless rules is supplied.
\end{document}
